\section{Discussion}

The reanalysis of the Tirosh melanoma single-cell RNA-seq dataset (3ca:20111) demonstrates that metabolic genes can separate cells into groups. This study does not establish universal, discrete malignant metabolic states. The analysis relies on a gene-membership snapshot of KEGG metabolic pathways, not on direct measurement of metabolic fluxes.

First, the clustering of all 4\,645 cells with $k=2$ yields a silhouette score of 0.377 and a mean pairwise adjusted Rand index (ARI) across seed fits of 0.99965, indicating strong reproducibility across seeds. However, the matched random gene set achieves a mean silhouette of 0.316 (SD 0.008) with an empirical $p$-value of 0.0476 from 20 exploratory runs. The 20 random gene sets provide a resolution of 1/21, and no multiple-testing correction was applied; thus the empirical $p$-value does not constitute proof of a biological state. The metabolic feature set's balanced accuracy (0.884) and macro-F1 (0.856) are comparable to matched random sets (balanced accuracy 0.896 $\pm$ 0.013; macro-F1 0.875 $\pm$ 0.015) and the highly variable gene baseline (balanced accuracy 0.892; macro-F1 0.897). No formal test of superiority was performed, and the results do not support a claim that metabolic genes define a distinct, biologically meaningful state.

Second, the strong correlation between PC1 and library complexity ($\rho = 0.650$, $p < 0.001$) suggests that technical variation contributes substantially to the observed separation. When expression features are residualized for complexity, the optimal $k$ shifts from 2 to 4, and the silhouette score drops to 0.224. This change indicates that the initial partition is sensitive to technical confounders. The adjusted Rand index remains high (0.99967), reflecting the stability of the partitioning algorithm rather than the biological validity of the groups. Residualization removes both technical and correlated biological signals; it cannot be used to assert that the initial separation was purely technical.

Third, the malignant-cell analysis reveals that clustering with $k=8$ yields a silhouette of 0.227 and seed ARI of 0.814, while joint patient and complexity residualization selects $k=2$ with silhouette 0.228 and seed ARI 0.970. The unadjusted optimum $k=8$ lies at the upper boundary of the candidate range (2--8), not at a biologically motivated value. The sample AMI for malignant cells is 0.742, indicating strong association with sample identity. The within-patient silhouette summaries (median 0.271, IQR 0.206--0.347) for the nine patients with at least 30 malignant cells do not constitute independent validation; they are derived from a single global PCA fit. The adjusted Rand index is not a measure of biological reproducibility but of algorithmic stability.

Fourth, the patient-blocked cross-validation assesses whether metabolic genes can predict known cell-type labels, not whether they define new states. The balanced accuracy and macro-F1 scores are high but are comparable to random and HVG baselines, indicating that the metabolic feature set does not outperform standard high-dimensional baselines.

\subsection{Limitations}

The dataset is a single cohort derived from the Tirosh melanoma study (GSE72056), and the Xiao2019 dataset is a differently processed subset of the same underlying data, precluding independent validation \cite{Tirosh2016,Xiao2019}. The analysis relies on a snapshot of KEGG metabolic gene membership (GMT) rather than direct measurement of metabolic fluxes. The detection filtering and highly variable gene selection were performed on the full dataset, not nested within training folds, which may introduce bias. PCA and scaling were applied within training folds, but the GroupKFold5 cross-validation uses sample groups to ensure no patient overlap between training and test folds; this is not a data leakage issue. K-means performs a forced partition of continuous data; the choice of $k$ does not prove the existence of discrete states. Complexity adjustment may remove genuine biological variation rather than eliminate technical artifact.

Future work should validate these findings in independent cohorts and, crucially, measure metabolic fluxes directly to determine whether transcriptional programs reflect functional metabolic states. The current results should be interpreted as exploratory, highlighting the need for rigorous validation before drawing biological conclusions.

\section{Conclusion}

Metabolic genes can separate single-cell melanoma cells into groups. This study does not establish universal, discrete malignant metabolic states. The analysis relies on a gene-membership snapshot of KEGG metabolic pathways, not on direct measurement of metabolic fluxes. The analysis is limited by the single-cohort design, the reliance on gene-membership snapshots, and the methodological constraints of unsupervised clustering. These limitations underscore the need for independent validation and direct metabolic measurements to establish whether transcriptional programs reflect functional metabolic states.
